% {{TITLE}}: notes with contents, theorems and proofs.
\documentclass[11pt]{article}
\usepackage[margin=1in]{geometry}
\usepackage{amsmath,amssymb,amsthm,mathtools}
\usepackage{graphicx}
\usepackage[colorlinks,allcolors=blue]{hyperref}
\usepackage{cleveref}

\theoremstyle{plain}
\newtheorem{theorem}{Theorem}[section]
\newtheorem{lemma}[theorem]{Lemma}
\newtheorem{proposition}[theorem]{Proposition}
\newtheorem{corollary}[theorem]{Corollary}
\theoremstyle{definition}
\newtheorem{definition}[theorem]{Definition}
\newtheorem{example}[theorem]{Example}
\theoremstyle{remark}
\newtheorem*{remark}{Remark}
\numberwithin{equation}{section}

\title{{{TITLE}}}
\author{{{AUTHOR}}}
\date{\today}

\begin{document}
\maketitle
\tableofcontents

\section{First topic}

\begin{definition}\label{def:hermitian}
An operator $A$ is \emph{Hermitian} if $A^\dagger = A$.
\end{definition}

\begin{theorem}\label{thm:real}
The eigenvalues of a Hermitian operator are real.
\end{theorem}
\begin{proof}
If $A\lvert\psi\rangle = a\lvert\psi\rangle$, then $a\langle\psi\vert\psi\rangle = \langle\psi\vert A\vert\psi\rangle
= \langle\psi\vert A^\dagger\vert\psi\rangle = a^*\langle\psi\vert\psi\rangle$, so $a = a^*$.
\end{proof}

\begin{remark}
\Cref{thm:real} uses only \cref{def:hermitian}.
\end{remark}

\section{Second topic}

\end{document}
